\documentclass[a4paper,12pt]{report}
\usepackage[utf8]{inputenc}
\usepackage[norsk]{babel}
\usepackage{hyperref}
\usepackage{geometry}
\geometry{margin=2.5cm}

\title{Norsk APA-manual\\ \large En nasjonal standard for norskspråklig APA-stil basert på APA 7th}
\author{Redaksjonen for norsk APA-stil}
\date{Versjon 1.12, 2025}

\begin{document}

\maketitle

\tableofcontents

\chapter{Innledning}
Denne manualen gir en oversikt over siteringsstilen APA i norsk språkform. APA-stilen er en siteringsstil som er utarbeidet av American Psychological Association.

\chapter{Forkortelser}
Tabellen viser eksempler på vanlige norske forkortelser som man bruker i APA.

\begin{tabular}{|l|l|}
\hline
\textbf{Begrep} & \textbf{Forkortelse} \\ \hline
Avsnitt & avsn. \\ \hline
Bind & bd. \\ \hline
Redaktør(er) & red. \\ \hline
Side(r) & s. \\ \hline
Uten årstall & u.å. \\ \hline
Utgave & utg. \\ \hline
\end{tabular}

\chapter{Henvisninger i teksten}
\section{Forholdet mellom henvisninger i tekst og oppsettet i referanselista}
APA-stilen baserer seg på «forfatter-årstall-systemet». Det skal være samsvar mellom kildehenvisningene i teksten og referanselista til slutt.

\section{Henvisning til en bestemt del av teksten}
Ved direkte sitat eller henvisning til spesifikk del, oppgi sidetall (s. 14), kapittel eller avsnitt.

\section{Sitat}
Korte sitater (under 40 ord) i teksten med anførselstegn.
Lange sitater (40 ord eller flere) som blokksitat med innrykk, uten anførselstegn.

\section{Parafrasering (indirekte sitat)}
Gjengi meningsinnholdet med egne ord. Kildehenvisning må inkluderes.

\section{Forfatternavn i henvisninger}
\begin{itemize}
    \item \textbf{1 eller 2 forfattere:} Nevn alle navnene hver gang (Hansen \& Olsen, 2020).
    \item \textbf{3 eller flere forfattere:} Bruk første forfatter et al. fra første henvisning (Hansen et al., 2020).
\end{itemize}

\section{Henvisning til kunstig intelligens}
Beskriv bruken av KI-verktøy i metodekapittelet eller i teksten. Referanser føres som programvare eller "Anonym"/Utgiver.

\chapter{Referanselista}
\section{Generelt}
Referanselista skal inneholde nok informasjon til å kunne identifisere og gjenfinne verket.
Fire hovedelementer: Forfatter, Publiseringsdato, Tittel, Kilde.

\section{Oppsett og sorteringsorden}
\begin{itemize}
    \item Alfabetisk rekkefølge etter første forfatters etternavn.
    \item Hengende innrykk.
    \item Dobbel linjeavstand (anbefalt).
\end{itemize}

\chapter{Eksempelsamling}

\section{Tekstlige verk}
\subsection{Artikler}
Forfatter, A. A. (År). Tittel på artikkel. \textit{Tidsskriftets navn, Volum}(Nummer), sidetall. DOI/URL

\subsection{Bøker}
Forfatter, A. A. (År). \textit{Tittel på boka} (Utg.). Forlag. DOI/URL

\section{Statistikk, datasett og programvare}
Forfatter/Organisasjon. (År). \textit{Tittel på datasett} [Datasett]. Utgiver. URL

\section{Audiovisuelt materiale}
Film, TV-serier, Podkast, Musikk, Bilder.

\section{Online media}
Nettsider, Sosiale medier.

\chapter{Endringslogg}
Oversikt over endringer fra tidligere versjoner.

\end{document}
